\documentclass[12pt,a4paper]{article}

% ── Packages ──────────────────────────────────────────────────────────────
\usepackage[margin=1in]{geometry}
\usepackage{graphicx}
\usepackage{booktabs}
\usepackage[colorlinks=true,linkcolor=blue!60!black,citecolor=blue!60!black,urlcolor=blue!60!black]{hyperref}
\usepackage[round,authoryear]{natbib}
\usepackage{setspace}
\usepackage{titlesec}
\usepackage{enumitem}
\usepackage[dvipsnames]{xcolor}

% ── Formatting ────────────────────────────────────────────────────────────
\onehalfspacing
\setlength{\parindent}{0pt}
\setlength{\parskip}{6pt}

\titleformat{\section}{\Large\bfseries}{\thesection.}{0.5em}{}
\titleformat{\subsection}{\large\bfseries}{\thesubsection}{0.5em}{}

% ══════════════════════════════════════════════════════════════════════════
\begin{document}

% ── Title Page ────────────────────────────────────────────────────────────
\begin{titlepage}
  \centering
  \vspace*{2cm}

  {\LARGE\bfseries Smart Campus Mobility:\\[6pt]
   A Pilot Study of Ride-Sharing Adoption\\[6pt]
   in Iraqi Universities\par}

  \vspace{1.5cm}
  {\Large Research Proposal\par}

  \vspace{2cm}
  {\large
    \textbf{Dr.\ Ali Al Ghanimi}\\[6pt]
    Assistant Professor\\[4pt]
    Department of Control Systems and Robotics\\[4pt]
    Faculty of Engineering, University of Kufa\\[4pt]
    Najaf, Iraq\\[8pt]
    \texttt{ali.alghanimi@uokufa.edu.iq}
  \par}

  \vfill
  {\large March 2026\par}
\end{titlepage}

% ── Table of Contents ─────────────────────────────────────────────────────
\newpage
\tableofcontents
\newpage

% ══════════════════════════════════════════════════════════════════════════
\section{Background and Motivation}
\label{sec:background}

Urban transportation in Iraqi cities faces persistent challenges: chronic traffic congestion, limited public transit infrastructure, rising fuel costs, and a lack of integrated mobility solutions. University campuses in Iraq, which serve tens of thousands of commuting students daily, represent a microcosm of these broader problems. Most students rely on private vehicles, informal taxis, or inconsistent minibus routes, resulting in long commute times, high costs, and significant environmental impact.

Ride-sharing platforms have transformed urban mobility worldwide. Empirical studies from the United States, Europe, China, and Southeast Asia have demonstrated that peer-to-peer ride-sharing can reduce vehicle kilometres traveled by 30--50\%, lower commuting costs by 20--40\%, and improve access to employment and education for underserved populations \citep{Shaheen2016,Tirachini2020,Mouratidis2021}. However, a critical gap persists in the literature: \textbf{there are zero peer-reviewed studies examining ride-sharing adoption in Iraq or comparable conflict-affected Middle Eastern contexts.} The socio-cultural factors that govern adoption in Iraq---including gender norms around mixed-gender travel, trust deficits in peer-to-peer transactions, cash-dominant economies, and security concerns---are fundamentally different from those studied in Western and East Asian settings.

\textbf{Toosila} is a fully developed, deployment-ready ride-sharing application specifically designed for Iraqi users. The platform incorporates culturally appropriate features such as gender-preference matching, family-member tracking, and integration with local payment norms. By piloting Toosila within the controlled environment of a university campus, this study aims to generate the first empirical evidence on ride-sharing adoption barriers and enablers in Iraq, contributing to both the transportation science literature and evidence-based mobility policy for Iraqi institutions.

% ══════════════════════════════════════════════════════════════════════════
\section{Research Questions}
\label{sec:rq}

This study addresses the following research questions:

\begin{description}[style=nextline, font=\bfseries]
  \item[RQ1] What are the primary barriers to ride-sharing adoption among Iraqi university students, and how do they differ from barriers reported in Western and East Asian contexts?
  \item[RQ2] How does interpersonal trust---both in co-riders and in the platform itself---evolve over the course of a six-week ride-sharing pilot, and what factors predict trust formation?
  \item[RQ3] To what extent do gender, faculty affiliation, and commute distance moderate willingness to adopt ride-sharing, and do gender-preference matching features mitigate adoption resistance?
  \item[RQ4] What are the measurable cost and time savings achieved by active ride-sharing participants compared to their pre-pilot commuting behaviour and compared to a waitlist control group?
\end{description}

% ══════════════════════════════════════════════════════════════════════════
\section{Methodology}
\label{sec:method}

\subsection{Study Design}

A quasi-experimental pre-post design with a waitlist control group will be employed. Participants will be randomly assigned to either the \textbf{treatment group} (immediate access to Toosila) or the \textbf{control group} (waitlisted for six weeks, then granted access). This design allows causal inference about the effects of ride-sharing access on commuting behaviour, cost, and attitudes.

\subsection{Participant Recruitment}

A minimum of \textbf{500 students} will be recruited from the Faculties of Engineering and Computer Science at the University of Kufa. Recruitment will proceed through:

\begin{itemize}[nosep]
  \item In-class announcements during regularly scheduled lectures (with instructor permission).
  \item QR-code posters placed in faculty buildings, cafeterias, and student common areas.
  \item Social media announcements on official university and departmental channels.
\end{itemize}

Eligibility criteria: (a)~currently enrolled undergraduate or graduate student, (b)~commutes to campus at least three days per week, (c)~owns a smartphone capable of running the Toosila application. Participants will provide informed consent and may withdraw at any time without academic penalty.

\subsection{Pre-Survey Instrument}

Prior to the pilot period, all participants will complete a structured pre-survey capturing the constructs shown in Table~\ref{tab:presurvey}.

\begin{table}[htbp]
  \centering
  \caption{Pre-survey constructs and measurement scales.}
  \label{tab:presurvey}
  \begin{tabular}{@{}lll@{}}
    \toprule
    \textbf{Construct} & \textbf{Items} & \textbf{Scale} \\
    \midrule
    Current commute mode & Single choice (car, taxi, minibus, walking, other) & Categorical \\
    Daily commute cost (IQD) & Open numeric & Continuous \\
    Daily commute time (min, one-way) & Open numeric & Continuous \\
    Willingness to share rides & 5 items & 7-point Likert \\
    Interpersonal trust & Adapted from \citet{Gefen2003}, 6 items & 7-point Likert \\
    Platform trust & Adapted from \citet{McKnight2002}, 5 items & 7-point Likert \\
    Gender comfort in shared transport & 4 items & 7-point Likert \\
    Technology readiness & Adapted from \citet{Parasuraman2015}, 8 items & 7-point Likert \\
    \bottomrule
  \end{tabular}
\end{table}

The survey will be administered in Arabic with back-translation validation.

\subsection{Six-Week Pilot Period}

Treatment-group participants will use the Toosila app for daily commuting over six weeks. The app will passively collect:

\begin{itemize}[nosep]
  \item Trip frequency, origin--destination pairs (anonymised to grid cells), and timestamps.
  \item Match rates (ride offers vs.\ ride requests fulfilled).
  \item In-app ratings and feedback.
  \item Feature usage logs (e.g., gender-preference filter activation, route deviation acceptance).
\end{itemize}

No GPS tracking data will be stored beyond anonymised grid-cell aggregation. Control-group participants will continue their normal commuting routines and complete weekly one-item commute-cost logs.

\subsection{Post-Survey Instrument}

After the six-week pilot, all participants will complete a post-survey including:

\begin{itemize}[nosep]
  \item Repeated measures from the pre-survey (commute cost, time, trust scales) to assess change.
  \item System Usability Scale (SUS) -- 10 items (treatment group only).
  \item Net Promoter Score (NPS) -- single item.
  \item Satisfaction with ride-sharing experience -- 6 items, 7-point Likert.
  \item Open-ended: perceived barriers, suggestions for improvement.
  \item Willingness to continue using ride-sharing post-study -- single item, 7-point Likert.
\end{itemize}

\subsection{Analysis Plan}

Quantitative analyses will include paired $t$-tests and mixed-effects models (group $\times$ time) for pre-post comparisons, structural equation modelling (SEM) for trust evolution pathways, and moderation analyses for gender and faculty effects. Qualitative responses will be analysed using thematic analysis \citep{Braun2006}. Effect sizes (Cohen's~$d$) and confidence intervals will be reported throughout.

% ══════════════════════════════════════════════════════════════════════════
\section{Data Collection and Privacy}
\label{sec:ethics}

All study procedures will be submitted for approval to the University of Kufa Research Ethics Committee prior to data collection. The following safeguards will be implemented:

\begin{description}[style=nextline, font=\bfseries]
  \item[Anonymisation] All participant data will be de-identified upon collection. A unique study ID will replace names and student numbers. The linking key will be stored on an encrypted, air-gapped device accessible only to the PI.
  \item[Location privacy] GPS coordinates will be aggregated to 500\,m grid cells before storage. Raw GPS traces will not be retained.
  \item[Data storage] Survey data will be stored on university-hosted servers. In-app behavioural data will be stored on encrypted cloud infrastructure with access restricted to the research team.
  \item[Informed consent] Participants will receive a bilingual (Arabic/English) consent form detailing data collection scope, storage duration (3~years post-publication), and their right to withdraw and request data deletion.
  \item[Compliance] The study will comply with Iraqi Ministry of Higher Education research guidelines and align with the principles of the Declaration of Helsinki.
\end{description}

% ══════════════════════════════════════════════════════════════════════════
\section{Expected Outcomes}
\label{sec:outcomes}

\begin{enumerate}[nosep]
  \item \textbf{First empirical dataset} on ride-sharing adoption in Iraq, filling a significant gap in the transportation literature for conflict-affected developing countries.
  \item \textbf{Identification of adoption barriers} specific to the Iraqi context (trust, gender norms, security, payment preferences), enabling culturally informed platform design.
  \item \textbf{Quantified cost and time savings} for ride-sharing participants, providing evidence for university administration to support campus mobility programmes.
  \item \textbf{Trust evolution model} demonstrating how repeated positive interactions change trust toward peer-to-peer mobility in low-trust environments.
  \item \textbf{Policy recommendations} for Iraqi universities considering subsidised or institutionally supported ride-sharing programmes.
\end{enumerate}

% ══════════════════════════════════════════════════════════════════════════
\section{Timeline}
\label{sec:timeline}

Table~\ref{tab:timeline} summarises the study schedule. The total active study duration is 12~weeks; manuscript preparation extends to week~20.

\begin{table}[htbp]
  \centering
  \caption{Study timeline.}
  \label{tab:timeline}
  \begin{tabular}{@{}llp{8cm}@{}}
    \toprule
    \textbf{Phase} & \textbf{Duration} & \textbf{Activities} \\
    \midrule
    Phase 1: Preparation       & Weeks 1--3   & Ethics approval, survey instrument validation, recruitment materials, app final testing \\
    Phase 2: Recruitment       & Weeks 4--5   & Participant enrolment, consent collection, pre-survey administration, randomisation \\
    Phase 3: Pilot Deployment  & Weeks 6--11  & Six-week active pilot (treatment group uses Toosila), weekly check-ins, technical support \\
    Phase 4: Post-Survey       & Week 12      & Post-survey administration, data export, initial data cleaning \\
    Phase 5: Analysis          & Weeks 13--20 & Statistical analysis, manuscript preparation, submission \\
    \bottomrule
  \end{tabular}
\end{table}

% ══════════════════════════════════════════════════════════════════════════
\section{Budget}
\label{sec:budget}

Since the Toosila application is already fully developed and deployed, the budget requirements for this pilot are minimal (Table~\ref{tab:budget}). No software development costs are required.

\begin{table}[htbp]
  \centering
  \caption{Estimated budget.}
  \label{tab:budget}
  \begin{tabular}{@{}lr@{}}
    \toprule
    \textbf{Item} & \textbf{Estimated Cost (USD)} \\
    \midrule
    Cloud hosting and server costs (3 months)                    & 150   \\
    Printed recruitment materials (posters, flyers)              & 100   \\
    Survey platform subscription (Arabic-compatible)             & 0     \\
    Participant incentives (mobile data top-ups for 500 students)& 1,000 \\
    Research assistant (1 RA, part-time, 12 weeks)               & 600   \\
    Contingency (10\%)                                           & 185   \\
    \midrule
    \textbf{Total}                                               & \textbf{2,035} \\
    \bottomrule
  \end{tabular}
\end{table}

% ══════════════════════════════════════════════════════════════════════════
\section{Target Journals for Publication}
\label{sec:journals}

The following journals are identified as suitable publication venues, listed by relevance and impact:

\begin{enumerate}[nosep]
  \item \textbf{Transportation Research Part A: Policy and Practice} (IF $\approx$ 6.3) -- ideal for adoption and policy findings.
  \item \textbf{Travel Behaviour and Society} (IF $\approx$ 5.2) -- focus on behavioural aspects of ride-sharing in developing countries.
  \item \textbf{Journal of Transport Geography} (IF $\approx$ 5.7) -- spatial analysis of campus ride-sharing patterns.
  \item \textbf{Sustainable Cities and Society} (IF $\approx$ 10.5) -- sustainability and smart campus mobility angle.
  \item \textbf{International Journal of Transportation Science and Technology} (IF $\approx$ 3.4) -- applied transportation technology, open access.
  \item \textbf{Case Studies on Transport Policy} (IF $\approx$ 2.8) -- suitable for a focused campus pilot case study if a shorter paper is preferred.
\end{enumerate}

A secondary publication targeting a regional conference (e.g., IEEE International Conference on Intelligent Transportation Systems) is also planned to disseminate preliminary results.

% ══════════════════════════════════════════════════════════════════════════
\section*{References}
\addcontentsline{toc}{section}{References}

\begin{thebibliography}{99}

\bibitem[Braun and Clarke(2006)]{Braun2006}
Braun, V., \& Clarke, V. (2006).
\newblock Using thematic analysis in psychology.
\newblock \textit{Qualitative Research in Psychology}, 3(2), 77--101.

\bibitem[Gefen et~al.(2003)]{Gefen2003}
Gefen, D., Karahanna, E., \& Straub, D.~W. (2003).
\newblock Trust and TAM in online shopping: An integrated model.
\newblock \textit{MIS Quarterly}, 27(1), 51--90.

\bibitem[McKnight et~al.(2002)]{McKnight2002}
McKnight, D.~H., Choudhury, V., \& Kacmar, C. (2002).
\newblock Developing and validating trust measures for e-commerce: An integrative typology.
\newblock \textit{Information Systems Research}, 13(3), 334--359.

\bibitem[Mouratidis et~al.(2021)]{Mouratidis2021}
Mouratidis, K., Peters, S., \& van Wee, B. (2021).
\newblock Transportation technologies, sharing economy, and teleactivities: Implications for built environment and travel.
\newblock \textit{Transportation Research Part D}, 92, 102716.

\bibitem[Parasuraman and Colby(2015)]{Parasuraman2015}
Parasuraman, A., \& Colby, C.~L. (2015).
\newblock An updated and streamlined technology readiness index: TRI~2.0.
\newblock \textit{Journal of Service Research}, 18(1), 59--74.

\bibitem[Shaheen et~al.(2016)]{Shaheen2016}
Shaheen, S., Chan, N., Bansal, A., \& Cohen, A. (2016).
\newblock Shared mobility: Definitions, industry developments, and early understanding.
\newblock \textit{UC Berkeley Transportation Sustainability Research Center}.

\bibitem[Tirachini(2020)]{Tirachini2020}
Tirachini, A. (2020).
\newblock Ride-hailing, travel behaviour and sustainable mobility: An international review.
\newblock \textit{Transportation}, 47, 2011--2047.

\end{thebibliography}

\end{document}
